%! Author = lili
%! Date = 7/9/26


\chapter{Tutorium 06.07.2026}
\untit{Präsenzaufgaben 2.5: Konfidenzintervalle \hfill 6.--10.\ Juli 2026}
\txo{Bearbeiten Sie auf jeden Fall Aufgabe 2.5.4. Falls Ihnen Zeit bleibt, bearbeiten Sie
zusätzlich auch Aufgabe 2.5.5.}

\section{PÜ 2.5.4}
Ein Politiker fällt wiederholt durch fragwürdige Äußerungen auf. Sein Wahlkampfteam
möchte wissen, inwieweit die Bevölkerung den Politiker weiterhin unterstützt.

In einer Stichprobe werden $200$ Personen befragt. $45$ Personen sprechen sich für den
Politiker aus.

\begin{auf}[Bestimmen Sie ein Konfidenzintervall zum Konfidenzniveau $0.975$ für die
Wahrscheinlichkeit, dass sich eine zufällig gewählte Person für den Politiker ausspricht.]
    NOCH KEINE LÖSUNG % TODO NOCH KEINE LÖSUNG
\end{auf}


\section{PÜ 2.5.5}
$n$-fach unabhängig wiederholten Bernoulli-Experiment) (6 Punkte)}
Sei $S_n$ die Anzahl der Erfolge in einem $n$-fach unabhängig wiederholten
Bernoulli-Experiment mit unbekannter Erfolgswahrscheinlichkeit $p \in (0,1)$. Wir
bezeichnen mit $J_n$ ein mit Hilfe von $S_n$ konstruiertes Konfidenzintervall für $p$.
Welche der folgenden Aussagen sind wahr? Begründen Sie Ihre Antwort.

\begin{auf}[Für jedes Konfidenzintervall $J_n$ für $p$ zur Konfidenzwahrscheinlichkeit
$1 - \alpha$ gilt $\P_p(p \in J_n) \ge \alpha$.]
% TODO OCR: Richtung/Größe der Ungleichung ggf. prüfen
NOCH KEINE LÖSUNG % TODO NOCH KEINE LÖSUNG
\end{auf}

\begin{auf}[Ein Konfidenzintervall zur Konfidenzwahrscheinlichkeit $1 - \alpha/2$ ist stets
doppelt so lang wie ein Konfidenzintervall zur Konfidenzwahrscheinlichkeit
$1 - \alpha$.]
NOCH KEINE LÖSUNG % TODO NOCH KEINE LÖSUNG
\end{auf}
